\documentclass[11pt]{article}
\usepackage[a4paper,margin=25mm]{geometry}
\usepackage[T1]{fontenc}
\usepackage[utf8]{inputenc}
\usepackage{lmodern}
\usepackage{amsmath,amssymb,amsthm}
\usepackage{hyperref}
\hypersetup{colorlinks=true,urlcolor=blue,linkcolor=blue,citecolor=blue}
\setlength{\emergencystretch}{3em}
\newtheorem{theorem}{Theorem}
\newtheorem{lemma}[theorem]{Lemma}
\newcommand{\crn}{\operatorname{cr}}
\newcommand{\N}{\mathbb{N}}
\title{An internal inconsistency in TLMC conjecture 00000001672}
\author{AlyciaBHZ, on behalf of the Omega Institute}
\date{5 October 2026}

\begin{document}
\maketitle

\begin{abstract}
The conjecture identifies $77$ as the Zarankiewicz formula's value at
$(7,7)$ and asserts both $\crn(K_{7,7})=77$ and that the formula holds
for $\min(m,n)\leq 7$. The formula's value is actually $81$.
Consequently the headline and its asserted restatement cannot both hold,
regardless of the true crossing number. We formalize this contradiction
for every natural-valued function in Lean~4, without defining crossing
numbers or using any axioms. Independently, Woodall's published
computer-assisted theorem gives $\crn(K_{7,7})=81$ and
$\crn(K_{7,9})=144$, refuting the isolated headline as well.
That crossing-number theorem is cited, not formalized here.
The description of $K_{7,7}$ as the smallest open instance is also false;
the smallest unsettled cases are $K_{7,11}$ and $K_{9,9}$.
\end{abstract}

\section{Statement and interpretation}
The English statement of conjecture \textbf{00000001672} is reproduced
verbatim below, with its mathematical notation typeset:
\begin{quote}
Definition: The Zarankiewicz formula asserts
$\crn(K_{m,n}) = \lfloor m/2\rfloor\allowbreak\lfloor(m-1)/2\rfloor
\allowbreak\lfloor n/2\rfloor\allowbreak\lfloor(n-1)/2\rfloor$.
Conjecture: $\crn(K_{7,7}) = 77$ (the value of the formula at the smallest
open instance); that is, the Zarankiewicz formula holds for
$\min(m,n) \leq 7$.
\end{quote}
The source is the repository file
\texttt{conjectures/00000001672.md}~\cite{tlmc}.

For positive integers $m,n$, let $K_{m,n}$ be the complete bipartite
graph with part sizes $m,n$. Its crossing number is the minimum number
of edge crossings in a plane drawing. Write the displayed product as
\[
Z(m,n)=\left\lfloor\frac m2\right\rfloor
\left\lfloor\frac{m-1}2\right\rfloor
\left\lfloor\frac n2\right\rfloor
\left\lfloor\frac{n-1}2\right\rfloor.
\]
The initial ``Definition'' specifies the formula under discussion; it
does not supply a proof that the formula holds for crossing numbers.

For a function $c:\N\times\N\longrightarrow\N$, distinguish the three
claims in the text:
\begin{align*}
H(c)&:\quad c(7,7)=77,\\
V&:\quad Z(7,7)=77,\\
F(c)&:\quad \forall m,n\in\N,\quad
\min(m,n)\leq7\ \Longrightarrow\ c(m,n)=Z(m,n).
\end{align*}
We read the declarative headline and the phrase ``that is'' as asserting
both $H(c)$ and $F(c)$; the parenthesis additionally asserts $V$.
This preserves the numerical headline, the exact product, the boundary
$7$, and the universal range of the restatement. We prove that
$H(c)\land F(c)$ is impossible, even after omitting $V$, and separately
that $V$ is false. Thus the entire collection of assertions is inconsistent.

Graph parameters are usually positive. The formalization also proves
the same impossibility when $F$ quantifies only over $m,n\geq1$.
Only $(7,7)$ is used, so the convention for zero-sized parts and the
natural subtraction at $m=0$ or $n=0$ cannot affect the contradiction.
For positive parameters, natural division by $2$ is precisely the floor
in the displayed expression.

\section{Full mathematical proof}
\begin{lemma}
$Z(7,7)=81$, and in particular $Z(7,7)\ne77$.
\end{lemma}
\begin{proof}
Both $\lfloor7/2\rfloor$ and $\lfloor(7-1)/2\rfloor
=\lfloor6/2\rfloor$ equal $3$. Therefore
\[
Z(7,7)=3\cdot3\cdot3\cdot3=81\ne77.
\]
\end{proof}

\begin{theorem}\label{thm:main}
For every $c:\N\times\N\to\N$,
\[
\neg\bigl(c(7,7)=77\ \land\
\forall m,n,\ \min(m,n)\leq7\Longrightarrow c(m,n)=Z(m,n)\bigr).
\]
The same conclusion holds if the universal quantifiers are restricted
to positive integers.
\end{theorem}
\begin{proof}
Suppose both assertions hold. Since $\min(7,7)=7\leq7$, the second
assertion gives $c(7,7)=Z(7,7)=81$. The first gives $c(7,7)=77$.
Transitivity yields $77=81$, a contradiction. The witness $(7,7)$ has
both parameters positive, which proves the positive-parameter variant
with exactly the same argument.
\end{proof}

In particular,
\[
F(c)\ \Longrightarrow\ c(7,7)\ne77.
\]
Even the weaker local assertion $c(7,7)=Z(7,7)$ contradicts the
headline. All these conclusions are independent of the definition or
actual value of the crossing number. Instantiating $c(m,n)$ with the
actual crossing number is a legitimate mathematical specialization of
the theorem; that specialization requires no crossing-number estimates.

\section{Scope and other readings}\label{sec:scope}
\textbf{Headline alone.} Woodall's computer-assisted theorem establishes
$\crn(K_{7,7})=81$, so the isolated headline $\crn(K_{7,7})=77$ is
false~\cite{woodall}. This is a cited mathematical refutation; Woodall's
crossing-number computation is not formalized in this submission.
The Lean proof independently refutes the headline and the asserted formula
range taken jointly, without depending on that computation.

\textbf{Formula range alone.} The assertion that the formula holds for
all $\min(m,n)\leq7$ is a separate claim. The Lean arithmetic argument
neither proves nor refutes that range for the actual crossing number.
Woodall settles the instances $(7,7)$ and $(7,9)$, not the full range:
$K_{7,11}$ remains among the smallest unsettled cases. Replacing $77$
by $81$ repairs the headline and agrees with Woodall's theorem, but does
not establish the universal range.

\textbf{Bare logical equivalence.} If ``that is'' is interpreted as the
single formula $H(c)\leftrightarrow F(c)$ without asserting either side,
the universal negation of this equivalence is false. Both sides can be
false. For example, let $c(m,n)=0$ for all $m,n$. Then $H(c)$ is false
because $0\ne77$, and $F(c)$ is false because at $(7,7)$ it would require
$0=81$. Hence $H(c)\leftrightarrow F(c)$ is true for this abstract
function. Lean formally exhibits this example; the zero function is
not proposed as the crossing-number function. For the actual crossing
number, Woodall makes $H$ false and the \emph{local} formula assertion
$\crn(K_{7,7})=Z(7,7)$ true, so the local equivalence is false.
This does not settle $F$, which asserts the formula for an entire range,
or the global bare equivalence $H\leftrightarrow F$.

\textbf{Literature context.} Kleitman's 1970 result establishes the
Zarankiewicz formula for complete bipartite graphs with smaller part at
most $6$~\cite{kleitman}. Woodall, in \emph{Cyclic-order graphs and
Zarankiewicz's crossing-number conjecture}, \emph{Journal of Graph Theory}
\textbf{17} (1993), 657--671, proved by a computer-assisted search
that~\cite{woodall}
\[
\crn(K_{7,7})=81,\qquad \crn(K_{7,9})=144
\]
Thus $K_{7,7}$ is a solved case, and both the isolated
headline and its description as the ``smallest open instance'' are false.
The smallest unsettled cases are $K_{7,11}$ and $K_{9,9}$.
These crossing-number results are cited, not formalized here, and are not
premises of the Lean proof. The formal result is the inconsistency of
the stated assertions, including $Z(7,7)=81\ne77$, for every function $c$.

\section{Lean statement and semantic correspondence}
The file \texttt{lean4/ZarankiewiczConsistency.lean} begins with the exact
statement and the interpretation formalized. It uses no imports, only
core Lean natural arithmetic. The symbol \texttt{Nat} is Lean's type of
natural numbers; the source supplies a local notation for that type.
Automatic implicit parameters are disabled.

The definition and main theorem, written using \texttt{Nat} and ASCII
notations equivalent to those in the source, are:
\begin{verbatim}
def Z (m n : Nat) : Nat :=
  (m / 2) * ((m - 1) / 2) * (n / 2) * ((n - 1) / 2)

theorem conjecture_inconsistent (cr : Nat -> Nat -> Nat) :
    Not (cr 7 7 = 77 /\
      forall m n, min m n <= 7 -> cr m n = Z m n)
\end{verbatim}
The fully qualified name is
\texttt{ZarankiewiczConsistency.conjecture\_inconsistent}.
The curried function \texttt{cr} represents $c$ above. It has no imposed
properties, so the result covers any natural-valued crossing-number
function. The conjunction faithfully retains the headline and the full
range assertion, rather than introducing a surrogate graph invariant.
The proof specializes the universal quantifiers at $7,7$.

The additional theorems are:
\begin{itemize}
\item \texttt{Z\_seven}: $Z(7,7)=81$.
\item \texttt{not\_value\_of\_formula}: $Z(7,7)\ne77$.
\item \texttt{local\_inconsistent}: the headline and the local formula
equality cannot both hold.
\item \texttt{formula\_range\_excludes\_77}: $F(c)\Rightarrow c(7,7)\ne77$.
\item \texttt{positive\_parameters\_inconsistent}: the main contradiction
with the additional premises $1\leq m$ and $1\leq n$ in the range assertion.
\item \texttt{equivalence\_reading\_has\_abstract\_model}: there exists a
natural-valued function satisfying the bare equivalence; the witness is
the zero function described in Section~\ref{sec:scope}.
\end{itemize}
No assertion about the actual crossing number is hidden in these auxiliary
theorems. Every claim in the proof of Theorem~\ref{thm:main}, including
the arithmetic evaluation, is covered by the formalization.

\section{Verification, axioms, and reproduction}
The pinned toolchain is \texttt{leanprover/lean4:v4.33.0}. The project
declares Mathlib \texttt{v4.33.0}; its manifest pins revision
\texttt{db584cd6d46c92f209a44c0f1c829460d327499d}. Although Mathlib is
declared in the project configuration, the proof needs only Lean's
core library. The project has one library root and one proof file.

Every theorem listed above ends with an axiom audit command. The actual
compiler output in \texttt{verification.txt} reports that all seven
theorems do not depend on any axioms, including:
\begin{verbatim}
'ZarankiewiczConsistency.conjecture_inconsistent'
  does not depend on any axioms
\end{verbatim}
The source contains no \texttt{sorry}, \texttt{admit},
\texttt{native\_decide}, or new axiom declarations. Closed arithmetic
goals use the ordinary kernel-checked \texttt{decide} tactic.
There are no unproved assumptions, graph-computation certificates, or
external numerical dependencies in the main theorem.

The final proof file compiled successfully with Lean \texttt{v4.33.0}.
The compiler output and exit status are saved in \texttt{verification.txt}.
A fresh project can be built with:
\begin{verbatim}
cd lean4
lake exe cache get && lake build
\end{verbatim}
The independent sanity check \texttt{verify.py} uses only Python's
standard library. It computes rational floors exactly, checks all four
factors, checks $Z(7,7)=81\ne77$, and checks that $(7,7)$ lies in the
specified range. It neither computes nor assumes a crossing number.
Its actual output is appended to \texttt{verification.txt}. To reproduce
the supplementary check and report from the submission directory:
\begin{verbatim}
python3 verify.py
pdflatex -interaction=nonstopmode -halt-on-error report.tex
pdflatex -interaction=nonstopmode -halt-on-error report.tex
\end{verbatim}

\paragraph{Attribution.}
Submitted by AlyciaBHZ on behalf of the Omega Institute (trureturing
project, \url{https://github.com/the-omega-institute/trureturing}).
No code from trureturing or another submission is vendored.

\begin{thebibliography}{9}
\bibitem{tlmc}
{\raggedright
The Last Math Competition, conjecture 00000001672, English statement.
\url{https://github.com/The-Last-Math-Competition/The-Last-Math-Competition/blob/main/conjectures/00000001672.md}.\par}
\bibitem{kleitman}
D.~J.~Kleitman, The crossing number of $K_{5,n}$,
\emph{Journal of Combinatorial Theory} \textbf{9}(4) (1970), 315--323.
\href{https://doi.org/10.1016/S0021-9800(70)80087-4}
{doi:10.1016/S0021-9800(70)80087-4}.
\bibitem{woodall}
D.~R.~Woodall, Cyclic-order graphs and Zarankiewicz's crossing-number
conjecture, \emph{Journal of Graph Theory} \textbf{17}(6) (1993), 657--671.
\href{https://doi.org/10.1002/jgt.3190170602}
{doi:10.1002/jgt.3190170602}.
\end{thebibliography}
\end{document}
